% problem-000398 8 电化学与电池
\section*{8. 电化学与电池}
（图：）

\subsection*{8-1}
计算在pH分别为4.0和6.0的条件下碳酸—碳酸盐体系中主要物种的浓度。

已知 $\mathrm { H _ { 2 } C O _ { 3 } }$ 的 $K _ { a _ { 1 } } = 4 . 5 \times 1 0 ^ { - 7 }$ $K _ { a _ { 2 } } = 4 . 7 \times 1 0 ^ { - 1 1 }$

\subsection*{8-2}
图中a和b分别是pH=4.4和6.1的两条直线，分别写出与a和b相对应的轴的物种发⽣转化的⽅程式。

\subsection*{8-3}
分别写出与直线c和d相对应的电极反应，说明其斜率为正或负的原因。

\subsection*{8-4}
在 $\mathrm { p H } = 4 . 0$ 的缓冲体系中加⼊ $\mathrm { U C l _ { 3 } } ,$ 写出反应⽅程式。

\subsection*{8-5}
在 $\mathrm { p H } = 8 . 0 { \sim } 1 2$ 之间，体系中 $\mathrm { U O _ { 2 } ( C O _ { 3 } ) _ { 3 } } ^ { 4 - }$ −和 $\mathrm { U } _ { 4 } \mathrm { O } _ { 9 } ( \mathrm { s } )$ 能否共存？说明理由； $\mathrm { U O _ { 2 } ( C O _ { 3 } ) _ { 3 } } ^ { 4 - }$ −和 $\mathrm { U O } _ { 2 } ( \mathsf { s } )$ 能否共存？说明理由。
